\documentclass[../main]{subfiles}
\begin{document}
\ifSubfilesClassLoaded{\setcounter{page}{19}}{}
\section{Коммунизм как реальность в общественном производстве}
\label{sec:04_kommunizm_realnost}

Утверждение о том, что коммунизм уже есть, может показаться слишком смелым. Во всяком случае, оно требует подтверждения. В глаза бросается обратное – отсутствие коммунизма, господство труда, ценность которого для того, кто трудится, определяется лишь в форме платы за труд – заработной платы. Может показаться, что клеточка (зародыш) коммунизма – непосредственно- коллективная форма свободной человеческой деятельности – нечто, не стоящее на магистральной дороге развития человеческого общества. Между тем, мы попытаемся показать, что это не так. Для этого недостаточно подобрать такие примеры, которые засвидетельствовали бы наличие коммунизма в урезанных рамках    эксплуататорского   классового    общества.    Нужно     рассмотреть непосредственно-коллективную        форму       свободной        деятельности, производящуюся в современном классовом обществе, со стороны возможностей её дальнейшего развёртывания.

Напомним, в своё время Ленин смог увидеть в таком свете сначала советы, а потом коммунистические субботники. «Прямо-таки гигантское значение в этом отношении имеет устройство рабочими, по их собственному почину, коммунистических субботников. Видимо, это только еще начало, но это начало необыкновенно большой важности. Это — начало переворота, более трудного, более существенного, более коренного, более решающего, чем свержение буржуазии, ибо это — победа над собственной косностью, распущенностью, мелкобуржуазным эгоизмом, над этими привычками, которые проклятый капитализм оставил в наследство рабочему и крестьянину. Когда эта победа будет закреплена, тогда и только тогда новая общественная дисциплина, социалистическая дисциплина будет создана, тогда и только тогда возврат назад, к капитализму, станет невозможным, коммунизм сделается действительно непобедимым». Ленин утверждал, что первый коммунистический субботник, устроенный 10 мая 1919 года железнодорожными рабочими Московско- Казанской железной дороги, имеет большее историческое значение, чем любая победа в Первой Мировой войне. Речь шла не просто о самостоятельном волеизъявлении рабочих, организовавших субботник, но и об существенном увеличении производительности труда в этом начинании, обусловленном тем, что для каждого работника это было его личным делом, где реализовалась личная заинтересованность в развитии железнодорожного хозяйства в интересах всего общества.

После публикации главы об основании коммунизма со стороны одного из наших читателей последовало важнейшее замечание: «Коммунистический способ производства подразумевает и производство материальных благ, а не просто пушкинский литературных кружок, способствующий становлению поэтов. Потому как "коммунистических отношений", основанных на эксплуатации для "избранных", сегодня полно» - утверждал читатель. Он потребовал предъявить свидетельства, что коммунизм уже есть, и каким он может быть именно в сфере материального производства. Правота читателя, на наш взгляд, тут заключается в двух аспектах. Во-первых, читатель выступил от имени пролетариата – людей, непосредственно занятых в материальном производстве путём продажи собственной рабочей силы, где коммунизм не наблюдается. Во- вторых, читатель ставит важнейший вопрос о коммунистическом материальном производстве, ведь именно оно определяет коммунизм в общественном масштабе, а не стремления, пожелания и даже действия политических революционеров.

Однако, коммунизм невозможен в чисто материальном производстве. Там, где человек непосредственно производит только вещи, а производство общественных отношений происходит в этом процессе совершенно независимо от него, от его целей, воли и сознания, ни о каком коммунизме не может идти речи, даже если само по себе производство вещей доставляет трудящемуся известное удовольствие от труда. В таком производстве человек является лишь материальной силой наряду с другими материальными силами. Он, в конечном счёте, выступает как вещь, воздействующая на другие вещи. Такой труд с точки зрения марксизма для наступления коммунизма не только может, но и должен быть уничтожен как человеческий труд и целиком передан машинам. Говоря об исторической роли капитала, Маркс специально отмечал важность создания прибавочного труда и универсальности человеческих потребностей как условия и предпосылки коммунизма. С этим капитал уже справляется и при капитализме. Но только этого недостаточно. Историческая роль капитала будет окончательно выполнена, а коммунистическое общество станет реальностью только тогда, «когда прекратится такой труд, при котором человек сам делает то, что он может заставить вещи делать для себя, для человека». В этом смысле человек не может действовать по-коммунистически, пока он не выведен за рамки чисто материального производства, пока он делает то, что могут делать вещи и, что важно, в этом процессе сам выступает как вещь. Поэтому пролетариат – единственный класс, который до конца последовательно заинтересован в том, чтобы историческая роль капитала была выполнена до конца. Это и есть уничтожение тех условий, в которых общество постоянно раскалывается на классы.

Автор замечания правильно отметил, что в сфере производства идеального коммунизма можно найти сколько угодно. Сюда можно отнести работу учёного и коллектива учёных, которые для того, чтобы дать людям научную истину, работают и день и ночь, да ещё и читают просветительские лекции в свободное время. Известный минимум материальных благ тут является условием, а не целью    ведения      коммунистической      деятельности.   Коммунистической деятельностью является работа писателя, художника, музыканта, которые трудятся для того, чтобы люди видели, слышали, чувствовали и понимали мир лучше, чтобы раздвинуть границы человеческих способностей, а не ради денег или славы и т.д. и т.п. Конечно, далеко не вся деятельность в сфере идеального является коммунистической. Но если указанные в предыдущей главе моменты непосредственно-коллективной формы свободной человеческой деятельности налицо, то мы имеем дело с локальным коммунизмом, несмотря на то, что коммунизм здесь существует в весьма урезанном виде.

Коммунизм здесь, как правильно нам укажут, далеко не полный. Во-первых, сами условия для него производятся за счёт чисто капиталистической эксплуатации человека человеком. Во-вторых, капитал, как господствующее общественное отношение, загоняет и здесь свободу в жёсткие рамки. Нельзя построить коммунизм, как способ материального производства, не выходя за пределы сферы идеального. И хотя коммунизм в сфере идеального на протяжении всей истории человечества играл важнейшую роль, так как в свободной деятельности рождались те идеи, которые не просто отражают, но и «творят» мир, всё же он не может играть решающего значения, так как не делает ещё общество коммунистическим. В-третьих, тут речь пока не идёт о преодолении общественного разделения труда, что является важнейшим условием для коммунизма. Универсальность человека здесь может быть только на уровне сознания, но не на уровне целостности человека как предметно- практического существа.

Однако, существует такое материальное производство, которое не является чисто материальным производством, а является производством человека в его целостности. Такое производство не просто окольным путём определяет, а включает в себя производство идеального – общественного сознания, общественного идеала, общественного целеполагания. Приведённый выше пример с коммунистическим субботником, в ходе которого рабочие занимались погрузкой и разгрузкой грузов, а также ремонтом подвижного состава, хорошо показывает это. Неприглядный, на первый взгляд, ручной труд здесь был не просто трудом по переносу тяжестей, а трудом по производству новой социальной реальности. Тут уместно вспомнить субботники по уборке сахарного тростника на Кубе: с помощью довольно примитивных орудий труда под палящим солнцем участники не просто рубили тростник, а добывали средства для модернизации производства, строительства школ, больниц – ради создания условий для развития всех представителей кубинского народа. Такой труд был и в колонии имени Горького и коммуне имени Дзержинского, которыми руководил Антон Семёнович Макаренко. Любое дело, даже выкапывание при помощи лопаты отхожих ям и возведение над ними деревянных домиков становилось здесь, в конечном счёте, производством нового человека, свободной человеческой деятельностью по меркам истины, добра и красоты.

Поэтому вопрос о свободной форме непосредственно-коллективной деятельности в материальном производстве нужно поставить так: какую связь с обществом, как целым в его развитии, такая деятельность воплощает и производит для того, кто трудится, и какое место и роль в развитии общественного целого она играет?

Нетрудно заметить, что приведённые выше примеры взяты из истории социализма. При социализме, самим фактом формального обобществления средств производства, снимается ряд препятствий для свободного изменения общественной жизни. Потому нам могут заметить, что эти примеры не совсем корректны. Но, во-первых, при снятии барьеров получает возможность реализации не то, чего вообще нет, и не то, для чего вообще нет условий, а то, для чего не было полного комплекса условий, и единственное, чего не хватало в этом комплексе – это упразднения господства капитала. То, что снятие барьеров даёт простор для развития коммунистических тенденций, означает, как минимум, что эти тенденции есть, что потенциал для становления такой деятельности кроется уже в капитализме.

Нам также могут возразить, что субботник – это что-то неустойчивое, ситуативное, не определяющее ход жизни вообще. Хотя в процессе субботника ставятся и достигаются общественные цели как личные цели участников субботника, за пределами субботника в материальном производстве действуют всё те же капиталистические законы: рабочая сила является товаром, а каждый отдельный пролетарий заключает договор купли-продажи рабочей силы со своим же классом, который выступает по отношению к нему как коллективный капиталист; закон стоимости сохраняет свою силу, если его действие специально не ограничивать, доводя тенденцию капитала к автоматизации производства до предела.

То есть, приведённое выше в качестве примера, само по себе ещё не делает общество коммунистическим. Замечания вполне верные, с этим мы спорить не станем. Но суть дела заключается не только и не столько в простой регистрации, простом фиксировании: дескать это – коммунизм, а это – не коммунизм, или не совсем коммунизм, а в тех тенденциях общества, потенциал которых заключается и виден в таких островках коммунизма. Именно в этом духе думал «ростках коммунизма» Ленин. Субботники его интересовали именно как ростки, зародыши. Потому он и обращает внимание, главным образом на две вещи: 1) свободное волеизъявление участников субботника по отношению и к средствам производства, и к самим себе; 2) значительный рост производительности труда во время субботников, достижение которого связано с этим свободным волеизъявлением.

Ещё одним существенным замечанием было бы указание на то, что в приведённых выше примерах деятельность осуществляется далеко не самыми передовыми даже для капитализма орудиями, и уж точно не на коммунистической технологической базе. Это напряжение человеческих сил происходит скорее от недостатка общественного развития: недостатка техники, недостатка   организации    технологических   процессов    даже    в  чисто капиталистическом плане. В этом, на наш взгляд, заключается самое резонное возражение против такого коммунизма, как против зародыша нового общества. Коммунизм по нужде, о котором говорилось выше, ещё не является коммунизмом по внутренней необходимости, а потому из него самого ещё не следует развитие коммунистического общества. Он является пусть и существенной, но только предпосылкой для такого общества. Здесь, опять-таки, нужно исследовать связь, преемственность со стороны способов построения коллективности между такого рода коммунизмом и коммунизмом, развитым на своём собственном основании, чему нужно посвятить отдельную работу.

Однако, отсылка к необходимости отдельной разработки этой проблемы вовсе не снимает с нас обязанности указать на практическую возможность этой связи в общих чертах, а также обязанности по предоставлению свидетельств существования непосредственно-коллективных форм свободной человеческой деятельности, основанных на передовых современных технологиях.

Связь между «коммунизмом по нужде» и коммунизмом на своём собственном базисе, их преемственность, может заключаться во-первых, в воспитании масс, в формировании привычки к свободному волеизъявлению по общественной необходимости в производстве общественной жизни; во-вторых, в создании коммунистических форм коллективной организации, которые в условиях преодоления общественной нужды, вызвавших их к жизни, вовсе не обязательно должны исчезнуть, а с учётом осуществляющегося в них воспитания в духе подлинной свободы, как осознанной общественной необходимости, стать формами коммунистической организации общественной жизни. Мы не утописты. Как именно это будет осуществляться, мы придумывать не станем. Это каждый раз зависит от сплетения конкретно-исторических обстоятельств. Мы только фиксируем возможность такого дальнейшего развития, заключающуюся в подобных ростках со стороны коллективно-организационной формы человеческой деятельности, предполагающих свободное и равное участие в ней.

Что же касается коммунизма, основанного на передовых технологиях, такие ростки мы видим в процессе совместного пользования современными информационными технологиями и создании их. Например, операционная система Linux и другие программы, созданные в рамках производства свободного ПО, созданы на коммунистических началах.

Оговоримся, что речь идёт не о всех программах с открытым кодом. Многие крупные современные проекты дали начало огромным корпорациям, и от того, что код открыт, деятельность по созданию программ не является свободной. И там, где это имеет место, корпорация полностью определяет политику проекта, спонсирует его и получает от этого выгоду. Код пишется программистами за плату. Коммунизма здесь нет (хотя есть возможность по превращению такой деятельности в коммунистическую). Так, например, пишутся вспомогательные продукты, нужные, нескольким компаниям сразу. Кооперация удешевляет разработку. Но, в конечном счёте, продукт тут производится для получения прибыли. Часто для вообще какого-либо использования такого продукта и, уж тем более, коммерческого, нужны большие ресурсы, имеющиеся только у крупных корпораций. У них обычно есть коммерческий сервис или софт, для которого разработанное является одной из условных «деталей», то есть условием для его работы и успешных продаж. Но корпорации монополизировали то, что технологически имеет серьёзный потенциал, выходя за рамки капитализма. Этот потенциал связан с технической возможностью как производства, так и пользования      программным       обеспечением,     благодаря       развитию информационных машин и сетей.

Началось свободное ПО, и часть проектов сейчас существует как проекты там, где нет товарно-денежных отношений. Более того, разработчиками ведётся борьба за то, чтобы из этого никто не извлекал коммерческой выгоды. Это и есть эмпирически данный коммунизм в рамках капитализма. Открытое программное обеспечение предполагает полный свободный доступ к программам и к кодам программ, а также предоставление в свободный доступ самих программ и их кодов, созданных с использованием полученных ранее свободных исходников. Например, операционная система Linux, оказалась конкурентоспособна по отношению к коммерческому Windows. Корпорации в итоге всячески пытаются подмять под себя и делают, продукты на основе открытого кода этой системы. Но сама она так и остаётся бесплатной.

Правда, созданное в рамках некоммерческого свободного ПО, далеко не всегда, работает так, как это задумывалось, на что обращают внимание пользователи. Но и коммерческое, на пути к стадии полностью готового продукта, проходит через это, а какие-то разработки вообще закрываются из-за их тупиковости. То же самое и здесь. Только тут все стадии открыты и видны всем. Недоделанное и тупиковое тоже оказывается достоянием общественности. Свободные программы – это вообще часто полуфабрикат, рассчитанный на приведение пользователем в состояние полной готовности. Это осложняет его конкурентоспособность в рамках капитализма в мире с общим низким уровнем информационной грамотности.

Корпорации развивают свой продукт с учётом этого уровня грамотности, доводя до возможности примитивно-интуитивного использования. Их продукт предназначен для людей, которым в условиях отчуждения легче заплатить, чтобы пользование было интуитивным, чем затратить время на то, чтобы разобраться в программе. Продукт корпораций не предполагает того, что пользователи создают новые программы или модифицируют существующие под себя путём дописывания или переписывания части кода. Использование их не требует особо высокой культуры и уж тем более понимания технологических процессов. Потому программы, произведённые крупными корпорациями, имеют намного более широкие возможности для сбыта. К тому же они часто более высокого качества, так как проработаны во всех деталях.

Свободное некоммерческое во всех смыслах ПО, напротив, требует знания хотя бы элементарных основ программирования, предполагает возможность, а часто и требует дописывания и переписывания. Оно заставляет пользователя самому становиться программистом хотя бы отчасти, и постоянно совершенствоваться в этом\footnote{Правда, это может портить нервы и не давать сосредоточится на выполняемой задаче. В итоге, вместо того чтобы делать дело, людям приходится чинить постоянно разваливающийся инструмент для Linux. Но нет никаких принципиальных препятствий (если смотреть со стороны потенциала) к преодолению этого и написания более качественных программ при возможности больших затрат общественного времени на их создание.}. Именно последнее обстоятельство, которое ограничивает чисто капиталистическую конкурентоспособность свободного ПО, говорит о его большем общественном потенциале, по отношению к коммерческому (речь не об отказе от удобного интерфейса, а о возможности самому создавать удобный интерфейс). В его производстве, способах передачи и использования заложен не просто потенциал технического развития, а потенциал развития человека – преодоления разделения труда, подтягивания всех участников процесса до уровня современной компьютерной грамотности. И хотя этот локальный коммунизм пытаются (и иногда успешно) подмять под себя капиталистические компании, его вполне можно отнести к росткам нового общества, пусть пока ещё слабым и нежным. Во-первых, он возник на базе самых передовых для капитализма технологий. Во-вторых, чётко очерчивает границу между капиталистическим и выходящим за рамки капитализма способом использования и создания этих технологий.

Важно ещё и то, что такое производство уже чисто технологически ведёт к уничтожению иерархии в производстве и раздельности пользователей и программистов. Оно не требует внешней связи между ними и специальных общественных «органов управления», обеспечивающих эту связь. Внешнее управление людьми, даже в самой чистой демократической форме – в форме решения вопросов путём общих голосований – здесь уже не требуется. В тенденции оно отмирает, становясь излишним.

Нам могут заметить, что, говоря о некоммерческом свободном производстве программного обеспечения, речь здесь опять идёт о сфере производства идеального. Действительно, процесс создания кода программ, как разновидности текста относится к идеальной деятельности, а код программы (но не сама программа в её обычной работе) является идеальным продуктом. Однако, это та идеальная деятельность, которая является важнейшей частью современного промышленного материального производства. Создание программного обеспечения – это труд по внедрению автоматизации в промышленное производство. И, хотя код программы является идеальным продуктом, функционирование современных промышленных машин на основе этих программ (то есть «работа» программ в промышленности) является чисто материальным машинным производством.

Конечно, в рамках производства свободного программного обеспечения производится множество бессмысленных и бесполезных программ. Это не так уж важно. Важно то, что производятся и могут производиться не только они, а вполне пригодные для использования в масштабах современной индустрии продукты. И производятся эти продукты коммунистически. Хотя часто эти проекты или не заканчиваются, или растягиваются на годы из-за того, что участники должны зарабатывать себе на жизнь, а потому времени на свободную деятельность у них не так много.

То же самое мы наблюдаем в процессе создания и использования современных баз данных. Опять-таки, далеко не все они производятся коммунистически. Но существуют большие общедоступные базы данных, предполагающие исключительно     совместное   создание    и    использование.    По-другому производиться они просто не могут. В противном случае они обессмысливались бы и становились бы непригодными для тех целей, ради которых и появились. Например, некоторые пространственные базы данных предполагают свободное участие в их формировании для создания образа пространства в интересах развития общества.

Тут тоже объективные закономерности человеческой деятельности в тенденции исключают обязательную иерархичность, характерную для капиталистического способа производства, включая в процесс всех участников как равных и свободных.

Структура базы данных тоже является идеальным продуктом, возразят нам. Но при создании баз данных идеальная деятельность настолько слита с материальным производством (в том числе и материальным отражением при помощи баз) и воспроизводством общественной жизни, что её при всей идеальности уж точно нельзя отнести к надстроечной деятельности. Как и в случае с ПО – это идеальный момент материального общественного производства, не абстрагированный (не в сознании, а реально) из этого производства в виде выделившихся форм общественного сознания.

Таким образом, мы здесь указываем не просто на идеальную деятельность, а на промышленное материальное производство, и на тенденции и потенции, заключённые в нём уже при капитализме, которые при снятии барьеров в виде частной собственности на средства производства могут развиться в коммунизм, составив как чисто технологическую, так и коллективно-организационную, воспитывающую, создающую нового человека, базу для него.

Дальнейшее развитие этой тенденции ведёт к автоматической фабрике в идее, в которой Маркс увидел прообраз будущего материального производства. Здесь также заложен потенциал реализации идей В. М. Глушкова об «общегосударственных» (всеобщественных), системах управления экономикой

(этой автоматической фабрикой в общественном масштабе). В перспективе (при условии снятия социальных барьеров) каждый может иметь доступ ко всей информации о состоянии всего общественного производства в реальном времени. Каждый может участвовать в управлении не просто производством вещей (производство вещей имеет смысл лишь постольку, поскольку оно служит главной цели производства), а производством человека – общественной жизни в её целостности.

Что интересно, участие во всех этих процессах предполагает производство и потребление идеального продукта всеми участниками, но такого, который не отчуждён и не может быть отчуждён от материального производства. В этом прослеживается возникновение нового типа общественного сознания, снимающая его отчуждённые формы.

И потребность в снятии социальных барьеров для свободной деятельности тоже создаётся современным капитализмом. Сам он вне зависимости от чьей- либо воли и сознания объективно ставит этот вопрос на повестку дня. Объективная логика движения современного капитала, постоянно порождающая войны и тем самым делает невозможными привычные способы и формы деятельности людей. Капитал вынуждает людей создавать качественно новые формы деятельности для решения проблем воспроизводства общественной жизни в противодействии разрушительной силе капитала.
\end{document}
